\def\student{Vic Segers}
\def\bedrijfspromotor{Tim Dupont}
\def\pxlpromotor{Tim Dupont}
\def\stagebedrijf{PXL Expertisecentrum Smart ICT}

\documentclass[a4paper]{article}
\usepackage{./PXLStageportfolio}

\begin{document}
  \pagenumbering{gobble}

  \begin{titlepage}
    \begin{figure}
      \includegraphics[width=0.5\linewidth]{images/logo_PXL_University_of_applied_sciences_and_arts.png}
    \end{figure}
    \begin{center}
      \vspace*{0.5cm}
      \Huge\textbf{\textcolor{pxlgreen}{Stageportfolio}}
      \par
      \vspace{0.5cm}
      \Large\textbf{\textcolor{pxlgreen}{PXL-IT}}
      \par
      \vspace{2.5cm}
      \normalsize\textbf{\stagebedrijf}
      \par
    \end{center}
    \vspace{2.5cm}
    \begin{tabbing}
      Student: ~~~~~~~~~~~~~~~ \= \textbf{\student}\\
      Bedrijfspromotor: \> \textbf{\bedrijfspromotor}\\
      Hogeschoolpromotor: \> \textbf{\pxlpromotor}
    \end{tabbing}
    \vfill
    \framebox{%
    \begin{minipage}{0.5\linewidth}
      Hogeschool PXL\\
      Elfde Liniestraat 24\\
      B-3500 Hasselt
      \bigbreak
      \hypersetup{linkcolor=blue}\href{https://www.pxl.be/}{www.pxl.be}
    \end{minipage}}
  \end{titlepage}

  \begingroup
  \hypersetup{linkcolor=black}
  \tableofcontents
  \endgroup

  \pagenumbering{arabic}

  \section{Informatiefiche}
    \begin{tabbing}
      Student: \textbf{\student}\\
      \\
      ~~~~~~ \= vic.segers10@gmail.com\\
      \> +32 476 62 57 66\\
      \> AON - AI \& Robotics
    \end{tabbing}
    \begin{tabbing}
      Bedrijf stageplek: \textbf{\stagebedrijf}\\
      \\
      ~~~~~~ \= Elfde-liniestraat 24\\
      \> B-3500 Hasselt
    \end{tabbing}
    \begin{tabbing}
      Bedrijfspromotor: \textbf{Dhr. \bedrijfspromotor}\\
      \\
      ~~~~~~ \= +32 495 46 55 25\\
      \> tim.dupont@pxl.be
    \end{tabbing}
    \begin{tabbing}
      Hogeschoolpromotor: \textbf{Dhr. \pxlpromotor}
    \end{tabbing}
    Stageproject: \textbf{Navigating Smart UAV systems in dynamic environments}
    \begin{addmargin}[5ex]{0pt}
      Het doel van de stage is om een flexibele multi-agentsysteem te ontwikkelen voor robots. Er zullen ook enkele demonstraties uitgewerkt moeten worden om dit systeem werkend te zien.
    \end{addmargin}
  \section{Plan van aanpak}
    \subsection{Situatieschets stagebedrijf + motivatie}
      % In de situatieschets geef je aan in welke omgeving het project zich afspeelt.
      % Omschrijf de aard van het stagebedrijf en de activiteiten van het stagebedrijf.
      % Wat is het doel van de organisatie? Welke producten of diensten worden er
      % geleverd? Welke zijn de Unique Selling Points van het bedrijf op vlak van IT
      % (in welk IT domein zijn ze gespecialiseerd, op welk vlak van de IT business
      % willen ze zich profileren, ...)? Wat is jouw motivering van het gekozen domein,
      % het gekozen project en het gekozen stagebedrijf?
      Het project zal door 3 stagestudenten ontwikkeld worden. Iedere student krijgt
      een specifieke taak en deze worden samen gebracht tot een geheel. De stage
      gaat door in het AI Lab op de Corda Campus.\par
      \stagebedrijf ontwikkelt en onderzoekt allerlei slimme toepassingen van de
      emerging techonologies en bundelt praktijkgericht onderzoek en dienstverlening
      op het snijvlak van IT en elektronica. Het spitst zich toe op mobiele technologie
      in de breedste zin van het woord en zoekt naar slimme oplossingen voor
      objecten en gemeenschappen: van slimme sensoren, uurwerken, telefoons en tablets,
      tot slimme gebouwen en smart cities. Daarbij staat de toepasbaarheid en
      inzetbaarheid van ontwikkelende toepassingen in de praktijk steeds voorop.\par
      Ik heb voor \stagebedrijf gekozen omdat ik er de kans kreeg om mijn IT-
      project verder uit te werken. Ook het kiezen van een meer onderzoekende stage
      leek mij niet slecht om na mijn bachelor door te schakelen naar een master.

      \newpage

    \subsection{Probleemstelling(en)}
      % Waarom er iets moet worden geproduceerd, blijkt uit de probleemstelling. Wat
      % is het probleem dat je moet gaan oplossen? Waarom werd het project opgestart,
      % in welke fase bevindt het project zich op het moment van de stage,...? Wat er
      % geproduceerd gaat worden, wordt in de doelstelling uiteengezet.
      Het project heeft al meerdere voorafgaande fasen gehad. Zo heeft een IT-project
      zich al bezig gehouden met een deel van dit onderwerp. En dit IT-project ging
      verder op reeds onderzochte informatie van een stage ook op \stagebedrijf.
      Hierdoor heeft het stageproject een basis verder op te bouwen.

    \newpage

    \subsection{Doelstelling(en)}
      % Door middel van een doelstelling wordt het gewenste eindresultaat van het project
      % omschreven. Onder omschrijving verstaan we een bondige en relevante
      % beschrijving van de doelstellingen die je wil verwezenlijken, de methodiek die je
      % kan/gaat aanwenden, ...
      Er moet een architectuur gemaakt worden voor multi-agentsystemen. Hiervoor
      worden een aantal voorbeelden geïmplementeerd worden ter demonstratie van
      de geleverde architectuur. Ook een hierarchical state machine zal ontwikkeld
      moeten worden voor om de computers beslissingen te laten nemen.\par
      Het begint bij een grondige analyse van verschillende multi-agent besluitvorming
      systemen. En een uitwerking van deze systemen. Met uiteindelijk ook een thesis.

    \newpage

    \subsection{Randvoorwaarden}
      % Welke afspraken heb je gemaakt met:
      % • het stagebedrijf, de bedrijfspromotor (werkregime, gedrag, …);
      % • de school, de hogeschoolpromotor (overlegfrequentie, manier van rapporteren, …);
      % • jezelf (vooropgestelde actiepunten, inzet, doel, …);
      % • de medestudenten indien je met meerdere studenten aan een project werkt.
      \begin{outline}
        \1 \textbf{Beslissingen:}
          \2[] De wijze waarop besluiten tot stand komen, bijvoorbeeld: je bedenkt een aantal
          mogelijke oplossingen met zijn voor- en nadelen en de bedrijfspromotor maakt hier een keuze uit.
        \1 \textbf{Beperkingen:}
          \2[] Omschrijf eventuele beperkingen waarmee je rekening moet houden. Dit kunnen
          beperkingen op materiaal, tijd, personeel, ruimte, kennis, e.d. zijn.
        \1 \textbf{Kritische succesfactoren:}
          \2[] Wat heb je zeker nodig om het project te laten slagen (bijvoorbeeld werkruimte, apparatuur,
          ondersteuning...). Dit kunnen eventueel ook middelen of werkruimten buiten het bedrijf zijn.
        \1 \textbf{Onzekerheden:}
          \2[] Het bereiken van de doelstellingen kan voor sommige projecten, afhankelijk van de beperkingen
          en de kritische succesfactoren, nog een vraagteken zijn. Vermeld hier deze onzekerheden.
        \1 \textbf{Afspraken:}
          \2 Studenten houden zich aan de glijdende werkuren van \stagebedrijf (I.e. 38u per week, starten $\leq$
          9 am, minimaal 30 min middagpauze)
          \2 Alles wordt ontwikkeld op eigen laptop of op de voorziene workstations, tijdens de werkuren pushen
          naar een met de begeleider afgesproken repository
            \3 Version Control verplicht te gebruiken voor Smart-ICT
            \3 Github repo te voorzien door begeleider
          \2 Voortgang bewaken door frequent stand-up meeting te doen
            \3 Dagelijks invullen in $\#$standup kanaal op Slack workspace vóór 09u15
            \3 Fysieke stand-up met alle stagairs $\rightarrow$ 1x per week op vrijdag
          \2 1x per sprint presentatie van resultaten naar de begeleider toe (sprint planning)
          \2 Insturingen
            \3 Portfolio’s en andere communicatie steeds via EPOS naar begeleider PXL 
            \3 Deadlines voor teksten op EPOS respecteren! 
            \3 \colorbox{yellow}{Geel markeren wat nieuw of aangepast is}
          \2 Solliciteren tijdens stageperiode
            \3 Flexibel
            \3 Afstemmen met begeleiders en team
            \3 Tijd inhalen
            \3 Andere afspraken uit PPT stagebegeleiding
          \2 Geen shenanigans
            \3 Geen speeltuin
            \3 Geen boksring
            \3 etc.
      \end{outline}

    \newpage

  \subsection{Tijdsplanning}
    Voeg hier je planning toe voor de 12 weken van je stageperiode.\\
    Gebruik hiervoor de technieken die je hebt aangeleerd in het opleidingsonderdeel
    Business Flow Advanced 1.\\
    Indien je stagebedrijf een andere methode hanteert voor het opstellen van een
    planning, dien je deze methode over te nemen.\par
    \underline{Effectief verloop van het project}\par
    Bijkomend moet op dezelfde wijze het effectieve verloop van het project
    weergegeven worden.\\
    Zo zal je waarschijnlijk ook merken dat de planning, door omstandigheden, niet voor
    honderd procent gevolgd kan worden. Maar een planning geeft jezelf en je promotoren
    wel een goed beeld van nog moet gebeuren: in welke fase je op een bepaald moment zit
    of zou moeten zitten. Dit kan heel eenvoudig door een bijkomende balk te gebruiken in
    een andere kleur.\par
    Wanneer jullie met meerderen aan een project werken, moet uit de planning
    duidelijk afgeleid kunnen worden welke taken voor welke student gelden.

  \subsection{Bronnen}
    Er dient altijd expliciet te worden aangegeven op welke bronnen en naslagwerken
    het eindwerk berust. Door systematisch hierin de bronnen te vermelden verliest
    men achteraf niet te veel tijd om de bronnen terug te gaan opzoeken. Het kan
    om schriftelijke, mondelinge een elektronische bronnen gaan.\par
    Voor de omschrijving gelden enkele vaste regels die vermeld staan in de
    cursus op BlackBoard.
  \section{Rapportage}
  \subsection{Wekelijkse rapportage}
  \subsection{Eindrapportage}
  \section{Terugkoppelingsformulieren}
  \subsection{Stagebespreking}
\end{document}